% GENERATED FILE. Edit canonical lessons, then run npm run generate:books.
% canonical-source-hash: fnv1a64:132732c31c69290b
% canonical-lessons: MW-C206-tijo, MW-C206-chautho, MW-C206-upar, MW-C206-par, MW-C206-algo

\chapter{The Third House, Upstairs}
\label{ch:mw-the-third-house-upstairs}
\hlchaptermodality{206}

\begin{chapteropening}
\textbf{By the end of this chapter:} \emph{I can say which house, which lane and which floor.}

Complete the last lesson of chapter 206: I can say which house, which lane and which floor.
\end{chapteropening}

\section[\texorpdfstring{\mw{तीजो} (tījo)}{तीजो (tījo)}]{\mw{तीजो} (tījo) — third}

\label{lesson:MW-C206-tijo}



\begin{quote}
Before the new one: say the Marwadi for to get stuck, then the Marwadi for a bus conductor.
\end{quote}

\subsection*{You'll want to know: \mw{तीजो}}
\textbf{\mw{तीजो}} — \emph{tījo} — \textquotedblleft{}third\textquotedblright{}.

\begin{quote}
\mw{तीजो} \mw{घर} \mw{म्हारो} \mw{है।}
\end{quote}

\begin{quote}
\emph{tījo ghar mhāro hai.}
\end{quote}

\begin{quote}
The third house is mine.
\end{quote}

\begin{cousinweb}
\textbf{\mw{तीजो}} (\emph{tījo}) is inherited from Sanskrit \emph{tṛtīya}, \textquotedblleft{}third\textquotedblright{}. CDIAL gives no Marwari form in this entry, but lists Hindi \emph{tījā}, \textquotedblleft{}third\textquotedblright{} (CDIAL 5912).
\end{cousinweb}

\begin{grammarlens}[title={What you've built}]
One more word for this chapter.
\end{grammarlens}

\subsection*{Your turn}
\begin{itemize}
\raggedright
  \item \emph{Say it:} \emph{tījo}
  \item \emph{Say it:} \emph{tījo}, once more
  \item \emph{Say it:} \emph{tījo ghar mhāro hai.} — The third house is mine.
  \item \emph{Recall:} say the Marwadi for to get stuck, then the Marwadi for a bus conductor, then say \emph{tījo} again
\end{itemize}

\begin{tcolorbox}[breakable,colback=teal!4,colframe=teal!35!black,title={Before you move on}]
What does \mw{तीजो} mean? (\textquotedblleft{}Third\textquotedblright{}.) Say it once more.
\end{tcolorbox}
\section[\texorpdfstring{\mw{चौथो} (chautho)}{चौथो (chautho)}]{\mw{चौथो} (chautho) — fourth}

\label{lesson:MW-C206-chautho}



\begin{quote}
Before the new one: say the Marwadi for a bus conductor, then the Marwadi for third.
\end{quote}

\subsection*{You'll want to know: \mw{चौथो}}
\textbf{\mw{चौथो}} — \emph{chautho} — \textquotedblleft{}fourth\textquotedblright{}.

\begin{quote}
\mw{मंदिर} \mw{चौथी} \mw{गळी} \mw{मांय} \mw{है।}
\end{quote}

\begin{quote}
\emph{mandir chauthī gaḷī mā̃y hai.}
\end{quote}

\begin{quote}
The temple is in the fourth lane.
\end{quote}

\begin{cousinweb}
\textbf{\mw{चौथो}} (\emph{chautho}) is inherited from Sanskrit \emph{caturtha}, \textquotedblleft{}fourth\textquotedblright{}; CDIAL lists Marwari \emph{cotho}, \textquotedblleft{}4th\textquotedblright{} (CDIAL 4600).
\end{cousinweb}

\begin{grammarlens}[title={What you've built}]
One more word for this chapter.
\end{grammarlens}

\subsection*{Your turn}
\begin{itemize}
\raggedright
  \item \emph{Say it:} \emph{chautho}
  \item \emph{Say it:} \emph{chautho}, once more
  \item \emph{Say it:} \emph{mandir chauthī gaḷī mā̃y hai.} — The temple is in the fourth lane.
  \item \emph{Recall:} say the Marwadi for a bus conductor, then the Marwadi for third, then say \emph{chautho} again
\end{itemize}

\begin{tcolorbox}[breakable,colback=teal!4,colframe=teal!35!black,title={Before you move on}]
What does \mw{चौथो} mean? (\textquotedblleft{}Fourth\textquotedblright{}.) Say it once more.
\end{tcolorbox}
\section[\texorpdfstring{\mw{ऊपर} (ūpar)}{ऊपर (ūpar)}]{\mw{ऊपर} (ūpar) — up; upstairs}

\label{lesson:MW-C206-upar}



\begin{quote}
Before the new one: say the Marwadi for third, then the Marwadi for fourth.
\end{quote}

\subsection*{You'll want to know: \mw{ऊपर}}
\textbf{\mw{ऊपर}} — \emph{ūpar} — \textquotedblleft{}up; upstairs\textquotedblright{}.

\begin{quote}
\mw{दफ्तर} \mw{ऊपर} \mw{है।}
\end{quote}

\begin{quote}
\emph{daphtar ūpar hai.}
\end{quote}

\begin{quote}
The office is upstairs.
\end{quote}

\begin{cousinweb}
\textbf{\mw{ऊपर}} (\emph{ūpar}) is inherited from a reconstructed \emph{*uppari}, \textquotedblleft{}above\textquotedblright{}; CDIAL lists Old Marwari \emph{ūpara}, \textquotedblleft{}on\textquotedblright{} (CDIAL 2333).
\end{cousinweb}

\begin{grammarlens}[title={What you've built}]
One more word for this chapter.
\end{grammarlens}

\subsection*{Your turn}
\begin{itemize}
\raggedright
  \item \emph{Say it:} \emph{ūpar}
  \item \emph{Say it:} \emph{ūpar}, once more
  \item \emph{Say it:} \emph{daphtar ūpar hai.} — The office is upstairs.
  \item \emph{Recall:} say the Marwadi for third, then the Marwadi for fourth, then say \emph{ūpar} again
\end{itemize}

\begin{tcolorbox}[breakable,colback=teal!4,colframe=teal!35!black,title={Before you move on}]
What does \mw{ऊपर} mean? (\textquotedblleft{}Up; upstairs\textquotedblright{}.) Say it once more.
\end{tcolorbox}
\section[\texorpdfstring{\mw{पार} (pār)}{पार (pār)}]{\mw{पार} (pār) — across; on the other side}

\label{lesson:MW-C206-par}



\begin{quote}
Before the new one: say the Marwadi for fourth, then the Marwadi for upstairs.
\end{quote}

\subsection*{You'll want to know: \mw{पार}}
\textbf{\mw{पार}} — \emph{pār} — \textquotedblleft{}across; on the other side\textquotedblright{}.

\begin{quote}
\mw{मंदिर} \mw{नदी} \mw{पार} \mw{है।}
\end{quote}

\begin{quote}
\emph{mandir nadī pār hai.}
\end{quote}

\begin{quote}
The temple is across the river.
\end{quote}

\begin{cousinweb}
\textbf{\mw{पार}} (\emph{pār}) is inherited from Sanskrit \emph{pāra}, \textquotedblleft{}the further bank\textquotedblright{}. CDIAL gives no Marwari form in this entry, but lists Hindi \emph{pār}, \textquotedblleft{}on the other side of, across\textquotedblright{} (CDIAL 8100).
\end{cousinweb}

\begin{grammarlens}[title={What you've built}]
One more word for this chapter.
\end{grammarlens}

\subsection*{Your turn}
\begin{itemize}
\raggedright
  \item \emph{Say it:} \emph{pār}
  \item \emph{Say it:} \emph{pār}, once more
  \item \emph{Say it:} \emph{mandir nadī pār hai.} — The temple is across the river.
  \item \emph{Recall:} say the Marwadi for fourth, then the Marwadi for upstairs, then say \emph{pār} again
\end{itemize}

\begin{tcolorbox}[breakable,colback=teal!4,colframe=teal!35!black,title={Before you move on}]
What does \mw{पार} mean? (\textquotedblleft{}Across; on the other side\textquotedblright{}.) Say it once more.
\end{tcolorbox}
\section[\texorpdfstring{\mw{अळगो} (aḷgo)}{अळगो (aḷgo)}]{\mw{अळगो} (aḷgo) — separate; apart}

\label{lesson:MW-C206-algo}



\begin{quote}
Before the new one: say the Marwadi for upstairs, then the Marwadi for on the other side.
\end{quote}

\subsection*{You'll want to know: \mw{अळगो}}
\textbf{\mw{अळगो}} — \emph{aḷgo} — \textquotedblleft{}separate; apart\textquotedblright{}.

\begin{quote}
\mw{म्हारो} \mw{घर} \mw{गांव} \mw{सूं} \mw{अळगो} \mw{है।}
\end{quote}

\begin{quote}
\emph{mhāro ghar gā̃v sū̃ aḷgo hai.}
\end{quote}

\begin{quote}
My house stands apart from the village.
\end{quote}

\begin{cousinweb}
\textbf{\mw{अळगो}} (\emph{aḷgo}) is inherited from Sanskrit \emph{alagna}, \textquotedblleft{}unconnected\textquotedblright{}; CDIAL lists Old Marwari \emph{alago}, \textquotedblleft{}separate\textquotedblright{} (CDIAL 700).
\end{cousinweb}

\begin{grammarlens}[title={What you've built}]
One more word for this chapter.
\end{grammarlens}

\subsection*{Your turn}
\begin{itemize}
\raggedright
  \item \emph{Say it:} \emph{aḷgo}
  \item \emph{Say it:} \emph{aḷgo}, once more
  \item \emph{Say it:} \emph{mhāro ghar gā̃v sū̃ aḷgo hai.} — My house stands apart from the village.
  \item \emph{Recall:} say the Marwadi for upstairs, then the Marwadi for on the other side, then say \emph{aḷgo} again
\end{itemize}

\begin{tcolorbox}[breakable,colback=teal!4,colframe=teal!35!black,title={Before you move on}]
What does \mw{अळगो} mean? (\textquotedblleft{}Separate; apart\textquotedblright{}.) Say it once more.
\end{tcolorbox}
